\section{Introdução}
\label{sec:introducao}

O avanço dos modelos de linguagem (\textit{Large Language Models} -- LLMs) transformou o paradigma de interação entre humanos e máquinas \textcolor{red}{\cite{mcpAtFirstGlance}}. No entanto, a ausência inicial de padrões de comunicação entre LLMs e ferramentas externas gerou uma fragmentação técnica que motivou o surgimento do \textit{Model Context Protocol} (MCP) \textcolor{red}{\cite{shi2026description}}. A adoção do protocolo tem crescido de forma acelerada, e um levantamento publicado no início de 2026 apontou que seus Kits de Desenvolvimento de Software (\textit{Software Development Kit} -- SDK) já registravam aproximadamente 97 milhões de \textit{downloads} mensais \cite{Avaya2026}. O principal diferencial do MCP reside no uso de ferramentas (\textit{tools}) descritas em linguagem natural, que o agente pode selecionar e invocar em tempo de execução conforme a tarefa a ser resolvida. Esse funcionamento impõe, contudo, uma dependência crítica da qualidade documental, uma vez que o modelo não tem acesso ao código que implementa as \textcolor{red}{ferramentas. As} decisões de usá-las, como invocá-las e quais parâmetros fornecer são \textcolor{red}{tomadas} com base em suas descrições \cite{shi2026description, 2025arXiv250814704L}. \textcolor{red}{Dessa} forma, assim como \textit{code smells} comprometem a manutenibilidade de \textit{softwares} tradicionais, descrições ambíguas ou incompletas comprometem o comportamento do agente, podendo induzir erros de planejamento, uso incorreto de parâmetros e alucinações \cite{2026arXiv260214878M}.

Estudos recentes investigam a prevalência de anomalias (ou \textit{smells}) em descrições de ferramentas MCP por meio de análise documental \textcolor{red}{\cite{2026arXiv260214878M}}. Essa abordagem avalia se a documentação comunica adequadamente o propósito, os parâmetros e as restrições de uso de uma ferramenta, examinando aspectos como clareza, completude e consistência, sem considerar a implementação técnica \cite{mcpAtFirstGlance}. Entretanto, evidências indicam que falhas relacionadas à configuração e à documentação representam 38,66\% dos problemas observados em servidores MCP \cite{Taraghi2026}. A elaboração e a manutenção dessas descrições dependem frequentemente de processos manuais desacoplados do código-fonte. \textcolor{red}{Cabe destacar que, no MCP, a descrição não determina nem gera a implementação da ferramenta, mas atua como um artefato de documentação redigido manualmente para um código já existente, sendo este o que define o comportamento efetivo da ferramenta em execução. A prevalência de inconsistências entre descrição e implementação observada em servidores reais \cite{shi2026description} evidencia que os dois artefatos podem divergir, o que motiva o uso do código-fonte como fonte adicional de evidência no processo de avaliação da qualidade da descrição.} Por esse motivo, existe o risco de surgirem discrepâncias semânticas entre a documentação e a implementação real. Diante disso, \textcolor{red}{existe uma lacuna na \textbf{análise da qualidade das descrições das ferramentas MCP, que, sem o contexto do código-fonte, pode não ser suficiente para a avaliação de sua qualidade}.}

A análise desse problema é relevante para a evolução sustentável de sistemas baseados em agentes de Inteligência Artificial (IA). Inconsistências estruturais podem reduzir a eficácia dos agentes e ampliar a superfície de exposição a vulnerabilidades, incluindo ataques de injeção de comandos \cite{Taraghi2026}. Além disso, investigações recentes identificaram a presença de anomalias em grande parte das ferramentas disponíveis no ecossistema MCP \cite{mcpAtFirstGlance}. Relatórios de segurança também alertam que ambiguidades documentais podem favorecer a execução de ações não autorizadas ou a exposição indevida de dados sensíveis \cite{NSA2026}. Nesse contexto, a incorporação de informações provenientes do código-fonte apresenta potencial para complementar abordagens tradicionais de avaliação de qualidade. Tal incorporação permite confrontar a descrição textual com a implementação efetiva da ferramenta, fornecendo evidências adicionais sobre possíveis inconsistências.

\textcolor{red}{Tendo em vista essa lacuna}, este trabalho tem como \textcolor{red}{objetivo geral \textbf{investigar se a inclusão do código-fonte da ferramenta influencia a avaliação da qualidade de descrições de ferramentas MCP.} Para alcançar esse objetivo, definem-se os seguintes objetivos específicos: (i) Avaliar as descrições de ferramentas MCP considerando o seu código; (ii) Avaliar se a adição da referência ao código da \textit{tool} altera o resultado da avaliação da descrição da \textit{tool}; e (iii) Identificar os atributos de avaliação das métricas afetados pela adição do código da \textit{tool} e realizar uma análise estatística sobre os resultados obtidos nas avaliações.}

Como parte da investigação, será conduzida uma análise sobre um conjunto de servidores MCP de código aberto, selecionados a partir de critérios de popularidade e maturidade do projeto. A avaliação irá se concentrar no confronto entre as descrições em linguagem natural das ferramentas e sua lógica de implementação, representada por informações estruturais extraídas do código-fonte. Espera-se que essa abordagem possibilite a identificação de inconsistências que não são facilmente observáveis por métodos exclusivamente textuais, como parâmetros implementados, mas não documentados, ou funcionalidades descritas sem correspondência evidente na implementação. Dessa forma, busca-se produzir evidências empíricas capazes de apoiar ou refutar a hipótese de que a incorporação do contexto técnico pode ampliar a capacidade de detecção de \textit{smells} em ferramentas MCP.

As próximas seções deste artigo estão organizadas da seguinte forma. A Seção~\ref{sec:fundamentacao} apresenta a Fundamentação Teórica sobre MCP, \textcolor{red}{qualidade no contexto de engenharia de software} e detecção de \textit{smells}. A Seção~\ref{sec:trabalhos_relacionados} discute os trabalhos relacionados. Por fim, a Seção~\ref{sec:materiais} descreve os materiais e métodos empregados na condução do estudo.
As próximas seções deste artigo estão organizadas da seguinte forma. A Seção~\ref{sec:fundamentacao} apresenta a Fundamentação Teórica sobre MCP, \textcolor{red}{qualidade no contexto de engenharia de software} e detecção de \textit{smells}. A Seção~\ref{sec:trabalhos_relacionados} discute os trabalhos relacionados. Por fim, a Seção~\ref{sec:materiais} descreve os materiais e métodos empregados na condução do estudo.